\section{Least classes, record primes, explicit families, and the carrier census}\label{sec:families}

\subsection{The least-class function and its record primes}\label{ssec:records}

The archive (§17.4 and §19, pp.~131 and 148--157) puts, for $i\ge1$, following Ionascu and Wilson \cite{IonascuWilson},
\[
\mathcal C_i=\Bigl\{n\ge2:\ \frac4n=\frac1x+\frac1y+\frac1z\ \text{with } x\le\frac{n+4i-1}4\Bigr\},\qquad h(n)=\min\{i:n\in\mathcal C_i\}.
\]
An odd prime $p$ is a \emph{strict record} if $h(q)<h(p)$ for every prime $q<p$.

\begin{lemma}\label{lem:leastclass}
Let $p\equiv1\pmod4$. The least denominator occurring in any solution at $p$ is the least occupied shell $a$. With $R=4a-p$ one has $h(p)=(R+1)/4$. For $p\equiv3\pmod4$, $h(p)=1$, and $h(2)=1$.
\end{lemma}
\status{Proved here.}
\begin{proof}
If $a$ is occupied, some solution has $a$ as a denominator, so the least denominator of that solution is at most $a$. Conversely, the least denominator $x$ of a solution is itself an occupied shell: it exceeds $p/4$ (Theorem~\ref{thm:shell}(a)) and is completed by the other two denominators. So the least occupied shell equals the least denominator over all solutions. A solution lies in $\mathcal C_i$ exactly when $4x-p\le4i-1$. For $p\equiv1\pmod4$ the residual $4x-p$ is $\equiv3\pmod4$, so the least $i$ is $(R+1)/4$. For $p\equiv3\pmod4$ the shell $(p+1)/4$ is occupied (Theorem~\ref{thm:shell}(c)). For $p=2$, take $4/2=1/1+1/2+1/2$ with $x=1\le5/4$.
\end{proof}

\begin{corollary}\label{cor:hle2}
If a prime $p$ does not lie in one of the six hard classes $1,121,169,289,361,529$ modulo $840$, then $h(p)\le2$. Hence every strict record after $73$ lies in a hard class.
\end{corollary}
\status{Proved here; found by the referee pass of version~1. It is the archive's Theorem~10.6 and Proposition~10.2 read through Lemma~\ref{lem:leastclass}. An independent computation of $h(p)$ for all $144{,}449{,}537$ primes below $3\cdot10^9$ finds no prime with $h(p)\ge3$ outside the hard classes.}
\begin{proof}
For $p=2$ and $p\equiv3\pmod4$, $h(p)=1$ by Lemma~\ref{lem:leastclass}. Let $p\equiv1\pmod4$.
\begin{itemize}
\item If $p\equiv5\pmod{12}$, the shell $(p+3)/4$ is occupied (Section~\ref{sec:core}, after Proposition~\ref{prop:r7}).
\item If $p\equiv13\pmod{24}$, the same shell is occupied (Proposition~\ref{prop:survivors}, first line of the proof).
\item If $p\equiv1\pmod{24}$ is not hard, Proposition~\ref{prop:survivors} shows that the shell with $R=3$ or the shell with $R=7$ is occupied.
\end{itemize}
In each case $R\le7$, so $h(p)\le2$. A strict record after $73$, where $h=2$, has $h\ge3$.
\end{proof}

\begin{proposition}[the record primes]\label{prop:records}
The strict records $p\le8{,}803{,}369$, with $h(p)$ and the residual $R$ of the least occupied shell, are exactly
\[
(2,1,2),\ (73,2,7),\ (1129,3,11),\ (1201,6,23),\ (21169,8,31),\ (118801,15,59),\ (8803369,27,107).
\]
For the odd records $R=4h(p)-1$; for $p=2$ the least denominator is $1$, so $R=4\cdot1-2=2$. Moreover $h(806521)=h(118801)=15$, so $806{,}521$ is not a strict record. There is no strict record between $8{,}803{,}369$ and $3\cdot10^9$. Below $3\cdot10^9$, no prime other than $8{,}803{,}369$ has $h\ge22$. The primes with $h=19$, $20$ or $21$ are
\begin{itemize}
\item $h=21$: $287{,}567{,}281$, $651{,}412{,}441$, $2{,}188{,}875{,}529$;
\item $h=20$: $778{,}103{,}881$, $1{,}749{,}233{,}641$, $2{,}535{,}613{,}921$, $2{,}809{,}257{,}889$;
\item $h=19$: $230{,}089{,}729$, $793{,}592{,}809$.
\end{itemize}
\end{proposition}
\status{Reproduced here: every prime up to $8{,}803{,}369$ was computed by Lemma~\ref{lem:leastclass} and Theorem~\ref{thm:shell}, using the residues of the divisors of $a^2$ modulo $R$ (script \note{esr\_records.py}, under three seconds).
\begin{itemize}
\item \emph{Source of the list.} Ionascu and Wilson print the record list \cite{IonascuWilson}. The archive (Proposition~17.6, pp.~131--132) recovers every printed record prime from the definition. It corrects one printed class label: $1201\in\mathcal C_6\setminus\mathcal C_5$, where the source prints $\mathcal C_4$. Its table (IW-RC) (p.~133) gives $R=2$ for $p=2$.
\item \emph{The range to $3\cdot10^9$.} It was found by the referee pass of version~1 with its own program, and recomputed here by an independent segmented program. The two agree on every record, on every prime with $h\ge19$, and on the histogram of $h$.
\item \emph{The shell at $p_*$.} At $p_*=8{,}803{,}369$ the least occupied shell has $R=107$, the shell of archive §3.
\end{itemize}}

For the five records whose residual $R$ is a safe prime ($R=2m+1$ with $m$ prime), and whose shell $a=(p+R)/4$ has exactly one non-residue prime factor modulo $R$, to exponent one, the archive defines a \emph{deficit} set $D\subset\ZZ/m$ from the discrete logarithms of the prime factors of $a$ (§19.1). It states the \emph{strict-record deficit-box conjecture} (§19.2): at such a record with non-empty deficit, $D$ is a difference-proper generalised arithmetic progression of rank at most two, and two induced \emph{odd sheets} cover $\ZZ/m$.
The deficits of the first three rows are empty. The two non-empty ones are proved to be a rank-two rectangle $\{11,14,15,18\}\subset\ZZ/29$ at $(118801,59)$ and the progression $\{23+42j:0\le j\le9\}\subset\ZZ/53$ at $(8803369,107)$.
The boundary example $p=806{,}521$ satisfies every hypothesis except strict recordness, and both conclusions fail there. So the hypothesis cannot be dropped.

\status{The conjecture is open, as the archive says. Its finite evidence is \textsf{reproduced} here from the definitions (script \note{esr\_deficits.py}, with the primitive root $g=2$ that the archive uses for $R=59$ and $R=107$):
\begin{itemize}
\item the deficits of the rows $p=73$, $1129$ and $1201$ are empty, and the row $p=21169$ is outside the domain because $m=15$ is composite;
\item at $p=118{,}801$: $a=29715=3\cdot5\cdot7\cdot283$, $D=\{11,14,15,18\}$, $D-D=\{3r+4s:|r|,|s|\le1\}$ with nine elements, offsets $u=11$ and $v=17$, and the two odd sheets cover $\ZZ/29$;
\item at $p=8{,}803{,}369$: $a=2{,}200{,}869=3^2\cdot11^2\cdot43\cdot47$, $D=\{23+42j:0\le j\le9\}$, $D-D=\{42k:|k|\le9\}$ with nineteen elements, offsets $29$ and $23$, and the sheets cover $\ZZ/53$; also the corollary of archive §16, that the middle target $-1$ is reached by exactly two exponent vectors and the exterior target not at all;
\item at $p=806{,}521$: $a=201{,}645=3^2\cdot5\cdot4481$, a deficit of $14$ points with $D-D=\ZZ/29$, and the odd sheets miss exactly $\{1,4,8,12,16,20,24,27\}$.
\end{itemize}
Another primitive root multiplies $D$ by a unit of $\ZZ/m$ (archive §19.1). The script confirms that the sizes of $D$ and $D-D$, and the number of missed points, do not change.
The failure of conclusion (i) at $806{,}521$ follows from $D-D=\ZZ/29$ by the archive's counting argument. A difference-proper progression of rank one with $14$ terms has exactly $27$ differences. A proper one of rank two has side lengths $\{2,7\}$, so a difference-proper one would have $3\cdot13=39>29$ distinct differences. That argument was read and is correct.
By Theorem~\ref{thm:shell}(b), the counts at $p_*=8{,}803{,}369$ ($|E_a|=0$, $|M_a|=2$) mean that exactly one unordered solution has least denominator $2{,}200{,}869$, namely the archive's (405) (p.~97), from the middle divisor $u=11^2$:
\[
\frac4{8803369}=\frac1{2200869}+\frac1{181085300330}+\frac1{3293760527702370}.
\]
The archive's Lean module covers the $R=107$ row only, as the archive states.}

\subsection{Fixed-divisor families, and four progressions on which the equation is solved}\label{ssec:families}

Theorem~\ref{thm:shell}(b) turns a single divisor into a solution. If the divisor is a fixed multiple of the cofactor $h=a/c$, the solution is given by polynomials in $n$ on a whole residue class. The archive states this for primes, as its constant-factor carriers (Lemma~24.218, p.~422). The form below holds for all $n$. It was found by the referee pass of version~1 and checked here.

\begin{proposition}[fixed-divisor families]\label{prop:fixeddivisor}
Let $R\ge3$ be odd, and let $c,s\ge1$ with $\gcd(c,R)=1$ and $s\mid c^2$. Let $n\ge1$ with $n\equiv-R\pmod{4c}$, and put $a=(n+R)/4$ and $h=a/c$, both integers.
\begin{enumerate}[label=(\roman*)]
\item \emph{Exterior.} If $n\equiv-cs^{-1}\pmod R$, then $4/n=1/a+1/y+1/z$ with the positive integers
\[
y=\frac{n(nsh+a)}R,\qquad z=\frac{c^2h/s+na}R .
\]
\item \emph{Middle.} If $n\equiv-4s\pmod R$, then $4/n=1/a+1/y+1/z$ with the positive integers
\[
y=\frac{n(a+s)}R,\qquad z=\frac{n(a+a^2/s)}R .
\]
\end{enumerate}
Each case is one residue class modulo $4cR$. If $\gcd(cR,35)=1$, that class meets all six hard classes modulo $840$ or none of them. It meets all six exactly when it is $\equiv1$ modulo $\gcd(4cR,24)$.
\end{proposition}
\begin{proof}
Put $S=na$. In both cases $n$ is a unit modulo $R$, because $\gcd(cs,R)=1$. So is $a$, because $4a\equiv n$. Hence $S$ is a unit modulo $R$. The proof of Theorem~\ref{thm:shell}(b) uses only this. A positive divisor $d$ of $S^2$ with $d\equiv-S\pmod R$ therefore gives positive integers $y=(d+S)/R$ and $z=(S^2/d+S)/R$ with $1/y+1/z=R/S$. Then $1/a+1/y+1/z=(n+R)/(na)=4/n$.
\begin{itemize}
\item In case (i), take $d=n^2u$ with $u=sh$. Here $u\mid c^2h^2=a^2$, and $d\equiv-S$ means $nu\equiv-a$, that is, $4u\equiv-1\pmod R$. That holds because $c(4sh+1)=s(n+R)+c\equiv sn+c\equiv0\pmod R$. Then $S^2/d=a^2/u=c^2h/s$.
\item In case (ii), take $d=nu$ with $u=s$. Here $s\mid c^2\mid a^2$, and $d\equiv-S$ means $s\equiv-a$, that is, $4s\equiv-n\pmod R$. Then $S^2/d=na^2/s$.
\end{itemize}
Since $R$ is odd and prime to $c$, the congruences modulo $4c$ and modulo $R$ define one class modulo $4cR$ (Chinese remainder theorem). If $\gcd(cR,35)=1$, then $\gcd(4cR,840)=\gcd(4cR,24)$. A class $r$ modulo $4cR$ meets a class $t$ modulo $840$ exactly when $r\equiv t$ modulo $\gcd(4cR,840)$. The six hard classes are all $\equiv1\pmod{24}$.
\end{proof}
\status{Proved here. For primes it is archive Lemma~24.218; the archive's proof of its Theorem~24.219 identifies the four families below as these middle and exterior channel maps.}

\begin{theorem}[four unconditional families]\label{thm:families}
For $k\ge0$ put $g_1=31k+30$, $g_2=31k+15$, $m_3=23k+15$, $m_4=23k+19$, $w_3=1116k+727$ and $w_4=4464k+3683$. Then $4/n_i=1/a_i+1/b_i+1/c_i$ with positive integers
\begin{center}\small
\begin{tabular}{@{}cllll@{}}
\toprule
$i$ & $n_i$ & $a_i$ & $b_i$ & $c_i$\\
\midrule
1 & $248k+209$ & $2g_1$ & $2n_1(k+1)$ & $2n_1g_1(k+1)$\\
2 & $248k+89$ & $2g_2$ & $n_2(2k+1)$ & $2n_2g_2(2k+1)$\\
3 & $17112k+11137$ & $186m_3$ & $12n_3m_3w_3$ & $124m_3w_3$\\
4 & $17112k+14113$ & $186m_4$ & $6n_4m_4w_4$ & $31m_4w_4$\\
\bottomrule
\end{tabular}
\end{center}
No primality of $n_i$ is needed.
\end{theorem}
\status{Proved here by direct expansion: $4a_ib_ic_i-n_i(b_ic_i+a_ic_i+a_ib_i)$ is the zero polynomial in $k$ for each $i$ (sympy), and the entries are positive for $k\ge0$; exact fractions agree for $k\le200$ (script \note{esr\_families.py}). Archive Theorem~24.219, p.~423.
\begin{itemize}
\item \emph{As fixed-divisor families.} The four families are instances of Proposition~\ref{prop:fixeddivisor}. With $(R,c,s)=(31,2,2)$ and $(31,1,1)$ in case (ii), and $(23,186,18)$ and $(23,186,36)$ in case (i), where $h=m_3$ and $h=m_4$, its formulas give exactly the table (checked by expansion). Two of the families are halves of larger instances (archive Proposition~24.222(ii)--(iii), p.~425):
\begin{itemize}
\item family 2 is half of the class $n\equiv89\pmod{124}$, which is $(31,1,1)$ itself;
\item family 3 is half of the class $n\equiv2581\pmod{8556}$, which is $(23,93,9)$ in case (i).
\end{itemize}
\item \emph{Rosati.} The archive shows that these triples are specialisations of Rosati's two 1954 parametrisations. It claims no independence from the classical parametrisations (archive abstract, p.~2, and §13.1). That comparison is located, not checked.
\end{itemize}}
\begin{proof}
All displayed linear factors have positive constant terms, so every entry is a positive integer for $k\ge0$. The identity $4a_ib_ic_i=n_i(b_ic_i+a_ic_i+a_ib_i)$ is a polynomial identity in $k$, verified by expanding both sides. Dividing by $n_ia_ib_ic_i$ gives the decomposition.
\end{proof}

Each progression meets all six hard classes modulo $840$. This is the last part of Proposition~\ref{prop:fixeddivisor}, and directly: since $\gcd(248,840)=8$ and $\gcd(17112,840)=24$, the numbers $n_1,n_2$ run through every class $\equiv1\pmod8$ modulo $840$, and $n_3,n_4$ through every class $\equiv1\pmod{24}$. Every hard class is $\equiv1\pmod{24}$, being the square of a unit.
Each $n_i(0)$ is prime to the step of its progression, so by Dirichlet's theorem every progression contains infinitely many primes in each hard class. Examples are the strict records $1201=n_1(4)\equiv361$ and $21169=n_2(85)\equiv169\pmod{840}$.
So the families solve the equation at infinitely many primes in every class that the reduction modulo $840$ leaves open (script \note{esr\_families.py} lists the least prime for each family and class).
This does not conflict with the obstruction to polynomial identities recalled in Section~\ref{sec:problem}. That obstruction is modulo $q$, and it applies only when $r$ is a square modulo $q$: here $11137\equiv5\pmod{23}$, and $5$ is not a square modulo $23$. The same holds for the other rows, since $14113\equiv14\pmod{23}$, $209\equiv23\pmod{31}$ and $89\equiv27\pmod{31}$ are non-squares. Proposition~\ref{prop:nosquare} below shows that this is forced: no fixed-divisor family covers a square class.

\subsection{The carrier census}\label{ssec:census}

The largest part of the archive's arithmetic is a sequence of \emph{sufficient-carrier censuses}. It occupies §18 and §24 from the base of the census (Theorem~24.100, p.~335) to the end of §24.3 (p.~477).
\begin{itemize}
\item \emph{Domain.} The census concerns only the primes $p\equiv25\pmod{72}$ whose shells $R=7$ and $R=11$ are both empty (archive Corollary~18.17 and Theorem~24.100). Such a prime lies in one of $30$ classes modulo $11088=16\cdot9\cdot7\cdot11$, the base rows $B_1\cup B_{25}$. The primes $p\equiv1$ and $p\equiv49\pmod{72}$ with these two shells empty are outside the census. Below $3\cdot10^6$ there are $2{,}316$ primes $p\equiv1\pmod{24}$ with both shells empty; $829$ are $\equiv1$, $682$ are $\equiv25$ and $805$ are $\equiv49\pmod{72}$ (computed here).
\item \emph{Rows.} The archive considers residue data of $p$ modulo a growing product of moduli. We call each such datum a \emph{row}. The moduli are $11088$, $5$, $19$, $23$, $31$ and $47$, then $71$, $191$, $83$, $167$ and $239$.
\item \emph{Forced rows.} A row is \emph{forced} when an explicit construction shows that, for every prime $p$ in the branch with that row, one of a finite list of shells is occupied. An example is the implication~(2272) on p.~431: for a prime $p\equiv25\pmod{72}$ whose shells $R=7$ and $R=11$ are empty and whose row lies in the domain but outside the survivor set, some shell with $R\in\{15,19,23,31,47\}$ is occupied.
\item \emph{Survivors.} The remaining rows are \emph{survivors}. Nothing is concluded for them.
\end{itemize}
The archive's front section (pp.~3--5) reports the following chain.
\begin{center}\footnotesize
\begin{tabular}{@{}>{\raggedright\arraybackslash}p{0.19\linewidth}rrr@{}}
\toprule
Step & Domain (rows) & Forced & Survivors\\
\midrule
six shells after the unit lift mod $71$ & $4{,}590{,}432{,}000$ & $4{,}212{,}008{,}640$ & $378{,}423{,}360$\\
$R=191$ & $71{,}900{,}438{,}400$ & $24{,}762{,}018{,}720$ & $47{,}138{,}419{,}680$\\
$Q=83$ & $3{,}865{,}350{,}413{,}760$ & $934{,}648{,}428{,}742$ & $2{,}930{,}701{,}985{,}018$\\
$R=167$ & $486{,}496{,}529{,}512{,}988$ & $115{,}851{,}891{,}806{,}170$ & $370{,}644{,}637{,}706{,}818$\\
$R=239$ & $88{,}213{,}423{,}774{,}222{,}684$ & $20{,}524{,}419{,}041{,}579{,}843$ & $67{,}689{,}004{,}732{,}642{,}841$\\
refresh of the $R=167$ coefficient & & $39{,}025{,}685{,}947{,}842$ & $67{,}649{,}979{,}046{,}694{,}999$\\
\bottomrule
\end{tabular}
\end{center}
Cumulatively $2{,}757{,}919{,}929{,}517{,}785{,}001$ of the $2{,}825{,}569{,}908{,}564{,}480{,000}$ rows of the final hard domain are forced, a proportion of $0.97606\ldots$.
The table is internally consistent (checked here):
\begin{itemize}
\item each domain is the previous survivor count times $q-1$, for the next modulus $q=191,83,167,239$;
\item the final domain is $4{,}590{,}432{,}000\cdot190\cdot82\cdot166\cdot238$;
\item forced plus survivors equals the domain in every row;
\item the fractions $69643/75900$ and $6257/75900$ of the first row are exact.
\end{itemize}

\status{Located. The arithmetic consistency of the reported numbers was checked here; the carrier theorems and the enumerations behind them were not re-run.
The archive states the following about its own checks.
\begin{itemize}
\item The enumerations are replayed by independent Python programs.
\item Bounded Lean modules check the literal census rows, totals, fractions, target residues and candidate orders, but not the affine and gcd reconstructions, carrier membership or the positive-denominator witnesses (pp.~3--5).
\item The ranking that selected $239$ after $167$ was corrected: on the full ten-coordinate object, $59$ ranks first and $239$ second. The $R=239$ step is kept as a valid non-greedy continuation, not as an optimal one (p.~4).
\end{itemize}
The archive calls these \emph{bounded sufficient-carrier censuses} and states that they are not shell-exhaustion, density, endpoint or optimality theorems (p.~2). A survivor row is only outside the carriers' images; it is not shown to be shell-empty.}

\paragraph{What the census reaches.}
Every carrier of the census is a fixed-divisor family in the sense of Proposition~\ref{prop:fixeddivisor}. On its row it fixes the residual $R$, a divisor $c$ of $a$ that the row determines, and a divisor $s$ of $c^2$, and uses $u=sh$ or $u=s$:
\begin{itemize}
\item for $R=15,19,23,31$ and for the factor-$31$ carrier at $R=23$, the maps listed in the proof of archive Theorem~24.228 (pp.~432--433), with $c=2,1,6$ or $12,2$ and $186$;
\item for $R=47$ and later, the sets $E_q(c)=\{-cs^{-1}\}$ and $M_q(c)=\{-4s\}$ over the divisors $s\mid c^2$, with $c=\gcd(a,N)$ read off the row (archive (2268)--(2269), p.~431, and the coefficients $c_Q$, pp.~462--463).
\end{itemize}
In each case $4cR$ divides the modulus of the row. The obstruction of Schinzel and Yamamoto therefore applies, in the following form.

\begin{proposition}[no fixed-divisor family covers a square class]\label{prop:nosquare}
In the setting of Proposition~\ref{prop:fixeddivisor}, the class modulo $4cR$ on which (i) or (ii) applies contains no square modulo $4cR$. So it contains no class that is a square modulo a multiple of $4cR$.
\end{proposition}
\begin{proof}
Suppose $n$ lies in the class and is a square modulo $4cR$.
\begin{itemize}
\item \emph{The prime $2$.} $n$ is odd, so $n\equiv1\pmod4$ and $R\equiv-n\equiv3\pmod4$. If $2\mid c$, then $8\mid4c$, so $n\equiv1\pmod8$, $R\equiv7\pmod8$ and $(2/R)=1$.
\item \emph{Odd primes $\ell\mid c$.} $-R\equiv n$ is a non-zero square modulo $\ell$, so $(-R/\ell)=1$. By reciprocity, $(\ell/R)=(R/\ell)(-1)^{(\ell-1)/2}=(-R/\ell)=1$, since $(R-1)/2$ is odd.
\end{itemize}
Hence $(c/R)=1$, and $(s/R)=1$ because $s\mid c^2$. But $n$ is a unit square modulo $R$, so $(n/R)=1$. In case (ii), $n\equiv-4s$ gives $(n/R)=(-1/R)(s/R)=-1$. In case (i), $n\equiv-cs^{-1}$ gives $(n/R)=(-1/R)(c/R)(s/R)=-1$. Either way this is a contradiction.
\end{proof}
\status{Proved here; found by the referee pass of version~1. It agrees with the archive's own residue sets (checked here):
\begin{itemize}
\item at $R=47$ the forced residues $F_{36}$ are exactly the $23$ non-residues modulo $47$, and $F_{18}$ consists of $21$ of them;
\item the unforced sets contain every quadratic residue of their modulus: $S_5=\{1,4\}$, $A_{25}=Q_{23}$, and $S_{19}$ and $S^\sharp_{31}$, whose excluded residues $15,18\bmod19$ and $15,23,27,29,30\bmod31$ are all non-residues.
\end{itemize}}

\begin{corollary}[the square rows of the census]\label{cor:censuslimit}
No row that is a square modulo the modulus of the census is ever forced by its carriers. All $30$ base rows are squares modulo $11088$. So the rows whose coordinates at the ten odd primes $5,19,23,31,47,71,191,83,167,239$ are all non-zero squares are never forced. They number
\[
\frac{2{,}825{,}569{,}908{,}564{,}480{,}000}{2^{10}}=2{,}759{,}345{,}613{,}832{,}500,
\]
that is, $1/1024$ of the final domain and $4.08\%$ of its $67{,}649{,}979{,}046{,}694{,}999$ survivors. Each further odd prime modulus added to the census halves this share, but no finite chain of such carriers removes it.
\end{corollary}
\begin{proof}
A forced row lies in a class of Proposition~\ref{prop:fixeddivisor} modulo $4cR$, and $4cR$ divides the modulus of the row. Proposition~\ref{prop:nosquare} applies. The base rows were checked directly: each is a square of a unit modulo $11088$ (script \note{esr\_referee\_checks.py}). The domain is $30\cdot\prod(q-1)$ over the ten odd primes, and exactly half of the units modulo each odd prime $q$ are squares.
\end{proof}

The survivors make up $2.39\%$ of the final domain. Each of the four lifts forced between $23\%$ and $35\%$ of its domain ($R=191$: $34.4\%$; $Q=83$: $24.2\%$; $R=167$: $23.8\%$; $R=239$: $23.3\%$). By Corollary~\ref{cor:censuslimit}, a residue census built from such carriers does not reach the square rows. Carriers that reach them are not fixed-divisor families; one kind of candidate uses the prime factors of $(p+R)/4$ (compare Proposition~\ref{prop:jacobi} and Theorem~I). The archive itself names the removal of the survivors as the open global obligation (p.~4).
